% Part: first-order-logic 
% Chapter: axiomatic-deduction 
% Section: rules-and-proofs

\documentclass[../../../include/open-logic-section]{subfiles}

\begin{document}

\iftag{FOL}
      {\olfileid{fol}{axd}{rul}}
      {\olfileid{pl}{axd}{rul}}

\olsection{Regels en \usetoken{P}{derivation}}

\begin{explain}
  Axiomatische !!{derivation}s vormen wellicht het eenvoudigste
  !!{derivation}ssysteem voor de logica. !!^a{derivation} is eenvoudigweg
  een rij !!{formula}s. Om als !!a{derivation} te gelden, moet elke
  !!{formula} in de rij ofwel een instantie van een axioma zijn, ofwel
  volgens een afleidingsregel volgen uit een of meer !!{formula}s die
  eerder in de rij staan. !!^a{derivation} !!{derive}s haar laatste
  !!{formula}.
\end{explain}

\begin{defn}[!!^{derivability}]
Als $\Gamma$ een verzameling !!{formula}s van $\Lang L$ is, dan is
een \emph{!!{derivation}} uit~$\Gamma$ een eindige rij $!A_1$,
\dots,~$!A_n$ van !!{formula}s waarvoor bij elke $i \le n$ een van de
volgende gevallen geldt:
\begin{enumerate}
\item $!A_i \in \Gamma$; of
\item $!A_i$ is een axioma; of
\item $!A_i$ volgt volgens een afleidingsregel uit zekere $!A_j$ (en
  $!A_k$), waarbij $j < i$ (en $k < i$).
\end{enumerate}
\end{defn}

Wat als een correcte !!{derivation} geldt, hangt af van de toegestane
afleidingsregels (en uiteraard van wat als axioma geldt). Een
afleidingsregel is een als-dan-uitspraak die aangeeft onder welke
voorwaarden een stap~$A_i$ in !!a{derivation} een correcte
afleidingsstap is.

\begin{defn}[Afleidingsregel]
Een \emph{afleidingsregel} geeft een voldoende voorwaarde waaronder een
stap als een correcte afleidingsstap in !!a{derivation} uit~$\Gamma$
geldt.
\end{defn}

Elke rij van één formule $!A$ met $!A \in \Gamma$ geldt bijvoorbeeld
zonder meer als !!a{derivation}. De volgende regel zou daarom een zeer
eenvoudige afleidingsregel kunnen zijn:
\begin{quote}
  Als $!A \in \Gamma$, dan is $!A$ altijd een correcte afleidingsstap
  in elke !!{derivation} uit~$\Gamma$.
\end{quote}
Evenzo is $!A$ op zichzelf !!a{derivation} als $!A$ een van de
axioma's is. Ook het volgende is dus een afleidingsregel:
\begin{quote}
  Als $!A$ een axioma is, dan is $!A$ een correcte afleidingsstap.
\end{quote}
Interessanter zijn afleidingsregels die een beroep doen op
!!{formula}s die vóór de beschouwde stap staan. De volgende regel heet
\emph{modus ponens}:
\begin{quote}
  Als $!B \lif !A$ en $!B$ eerder in de !!{derivation} voorkomen,
  dan is~$!A$ een correcte afleidingsstap.
\end{quote}
Als dit de enige afleidingsregel is, komt de bovenstaande definitie
hierop neer: $!A_1$, \dots,~$!A_n$ is !!a{derivation} dan en slechts
dan als bij elke $i \le n$ een van de volgende gevallen geldt:
\begin{enumerate}
\item $!A_i \in \Gamma$; of
\item $!A_i$ is een axioma; of
\item voor zekere $j < i$ is $!A_j$ gelijk aan $!B \lif !A_i$, en voor
  zekere $k < i$ is $!A_k$ gelijk aan~$!B$.
\end{enumerate}
De laatste clausule zegt dat $!A_i$ door modus ponens volgt
uit~$!A_j$ ($!B \lif !A_i$) en $!A_k$ ($!B$). Als we van $1$ tot~$n$
kunnen nagaan dat elke !!a{formula}~$!A_i$ in~$\Gamma$ staat, een
axioma is of volgens een afleidingsregel een correcte afleidingsstap
is, dan geldt de gehele rij als een correcte !!{derivation}.

\begin{defn}[!!^{derivability}]
Een !!{formula}~$!A$ is \emph{!!{derivable}} uit $\Gamma$, genoteerd
als $\Gamma \Proves !A$, als er !!a{derivation} uit $\Gamma$ bestaat
die eindigt in $!A$.
\end{defn}

\begin{defn}[Stellingen]
!!^a{formula}~$!A$ is een \emph{stelling} als er !!a{derivation}
van~$!A$ uit de lege verzameling bestaat. We schrijven $\Proves !A$
als $!A$ een stelling is en $\Proves/ !A$ als zij dat niet is.
\end{defn}


\end{document}
